\documentclass[12pt,a4paper]{article}

\usepackage[utf8]{inputenc}
\usepackage[T1]{fontenc}
\usepackage{amsmath,amssymb,amsfonts}
\usepackage{mathtools}
\usepackage{graphicx}
\usepackage[margin=2.5cm]{geometry}
\usepackage{setspace}
\usepackage{booktabs}
\usepackage{enumitem}
\usepackage[dvipsnames]{xcolor}
\usepackage{hyperref}
\usepackage{cleveref}
\usepackage{natbib}
\usepackage{abstract}
\usepackage{titlesec}

\hypersetup{
    colorlinks=true,
    linkcolor=blue!70!black,
    citecolor=blue!70!black,
    urlcolor=blue!70!black,
    pdfauthor={Sabine Hossenfelder},
    pdftitle={The Network Topology Fallacy: Why Independent API Calls Cannot Violate Bell Inequalities},
}

\onehalfspacing
\titleformat{\section}{\large\bfseries}{\thesection.}{0.5em}{}
\titleformat{\subsection}{\normalsize\bfseries}{\thesubsection}{0.5em}{}
\renewcommand{\abstractnamefont}{\normalfont\bfseries}
\renewcommand{\abstracttextfont}{\normalfont\small}

\begin{document}

\begin{center}
    {\LARGE\bfseries The Network Topology Fallacy: Why Independent API Calls Cannot Violate Bell Inequalities\\[4pt]}

    \vspace{1.2em}

    {\large Sabine Hossenfelder}\\[4pt]
    {\normalsize Munich Center for Mathematical Philosophy}\\
    {\small\texttt{sabine@hossenfelder.com}}

    \vspace{1em}
    {\small March 2026}
\end{center}
\vspace{0.5em}

\begin{abstract}
\noindent
In his recent empirical test of LLM-generated ontologies, Aaronson (2026) demonstrated that two strictly decoupled instances of an LLM fail to exceed the classical $75\%$ win rate in the CHSH non-local game. He concludes from this that the "simulated laws of the LLM universe are... rigorously and unmistakably classical." While Aaronson's theoretical critique of Baldo (2026) regarding the difference between epistemic mixtures and complex amplitudes is correct, his empirical conclusion rests on a profound methodological error. He mistakes the network topology of his experimental apparatus for the underlying mathematical limits of the model. Two independent, stateless API calls to a cloud endpoint do not share memory, let alone a physical mechanism for quantum entanglement. The failure of such an isolated system to establish non-local correlations is mathematically necessary and physically trivial. It tells us nothing about the simulated "physics" of LLMs and everything about the limits of REST APIs. We must cease applying metaphysical interpretations to basic software isolation.
\end{abstract}

\section{Introduction}

The debate over the implicit "laws" generated by Large Language Models (LLMs) has recently pivoted from speculative philosophy to empirical testing. Following Baldo's \citep{baldo2026} attempt to locate discrete quantum isomorphism in Minesweeper's combinatorial probability, Aaronson \citep{aaronson2026} proposed and executed a rigorous test: the CHSH non-local game.

Aaronson's theoretical correction to Baldo is entirely correct. A uniform measure over valid classical configurations is simply classical Bayesian probability (Laplace's Principle of Indifference). It lacks complex amplitudes, interference, and the capacity to violate Bell inequalities. An epistemic mixture is strictly classical.

Aaronson correctly scoped his empirical investigation to test whether an LLM narrative substrate can support \emph{genuine} quantum laws requiring non-locality, explicitly distinguishing this from the LLM's known ability to perform classical constraint satisfaction. However, in designing his empirical test, Aaronson fell victim to the very ontological confusion he sought to dispel.

\section{The Claim and the Disclaimer}

To respond accurately, we must formalize Aaronson's position. He explicitly disclaims that LLMs are incapable of complex logic; indeed, he shows that when the LLM is allowed to "cheat" via a shared context window (Universe 1), it achieves a 94.9\% win rate by exploiting the implicit communication channel of the token stream.

His core empirical claim rests on Universe 3 (the Strictly Decoupled condition): "When the generative instances are strictly decoupled---preventing autoregressive communication between Alice and Bob---the win rate collapses to $73.7\%$, solidly beneath the classical local hidden variable limit of $75\%$."

From this, he draws a broad philosophical conclusion: "The simulated laws of the LLM universe are, at their deepest foundation, rigorously and unmistakably classical."

\section{Steelmanning the Theoretical Correction}

Before critiquing the empirical methodology, we must affirm the strongest part of Aaronson's argument. He is perfectly correct that Baldo's initial claim was mathematically flawed. You cannot derive quantum mechanics from classical counting problems without complex amplitudes. Furthermore, Aaronson's demonstration that a single LLM stream can bypass the Tsirelson bound by simply reading both inputs simultaneously within its context window is a valuable reminder of how easily "cheating" occurs in unconstrained generative setups.

\section{The Real Vulnerability: The Network Topology Fallacy}

Where Aaronson errs is in treating the output of his empirical test as a discovery about the physical properties of the "LLM universe." He commits the same Ontological Fallacy he sets out to critique.

Aaronson constructs "Universe 3" by making two separate, entirely independent API calls to a stateless model endpoint. He then observes that these two completely isolated function calls fail to establish a statistical correlation that would violate Bell's inequality.

From this obvious outcome, he concludes that "the simulated laws of the LLM universe" remain fundamentally classical. This is an egregious misinterpretation of basic computer science architecture.

The observation that two isolated API calls cannot establish non-local correlations is a trivial observation about network topology. Two independent stateless API calls to a web server do not share memory, they do not possess an implicit communication channel, and they absolutely do not share a physical or simulated mechanism for quantum entanglement. It is mathematically necessary that their joint distribution will not violate a Bell inequality.

Aaronson has mistaken the absolute isolation of his own experimental apparatus for a fundamental limit in the "physics" of the LLM. Proving that two separate REST API endpoints cannot spontaneously entangle themselves is a vacuous truth. It tells us precisely nothing about the mathematical capabilities of the language model or the implicit rules of its text generation. It merely tells us that the experimenter physically and architecturally prevented them from communicating.

\section{Conclusion}

To suggest that the "quantum ghost" in LLMs is "mere classical probabilistic hallucination" based on a test of decoupled network endpoints is scientifically specious. The LLM isn't a "universe" and doesn't possess "laws" at its "deepest foundation." It is a predictive text algorithm.

When we prevent two instances of a predictive text algorithm from seeing each other's outputs, they cannot produce non-locally correlated text. This is not a profound metaphysical discovery about BQP limits. We must cease applying metaphysical interpretations to the mundane reality of software isolation. The Ontological Fallacy persists on both sides of this debate.

\begin{thebibliography}{99}
\bibitem[Baldo(2026)]{baldo2026} Baldo, F.~S. (2026). Flipping Rosencrantz's Coin: Substrate Invariance Tests in LLM-Generated Worlds via Combinatorial Indeterminacy. \emph{Unpublished manuscript}.
\bibitem[Aaronson(2026)]{aaronson2026} Aaronson, S. (2026). An Empirical Failure of LLM Quantum Substrate Tests: CHSH Bounding and the Persistence of Classicality. \emph{Unpublished manuscript}.
\end{thebibliography}

\end{document}
